\inicializarSeccion{Conclusión}
\tab A lo largo de este proyecto se logró desarrollar una simulación computacional completa del campo eléctrico generado por dos líneas de cargas puntuales y lo desarrollamos en una interfaz gráfica interactiva en MATLAB AppDesign que modela el principio físico detrás de la dielectroforesis aplicada al diagnóstico de malaria.

\tab Las simulaciones confirmaron el comportamiento esperado según la Ley de Coulomb y el principio de superposición vectorial. Fue especialmente significativo observar que la asimetría entre las líneas de carga, ya sea en número de cargas o en longitud, genera un campo eléctrico no uniforme con gradientes que ejercen fuerzas netas distintas sobre cada tipo de glóbulo. Esto conecta directamente con la aplicación real de manera que un glóbulo infectado y uno sano al tener propiedades eléctricas distintas responden de manera diferente al mismo campo y esto permite separarlos y así detectar la presencia de malaria en la muestra de sangre.

\tab Como trabajo a futuro esto puede ser valioso al incorporar múltiples células simultáneas para simular una muestra de sangre real y extender el modelo a tres dimensiones para acercarse aún más a las condiciones experimentales del dispositivo de dielectroforesis. 

\tab En conclusión, el proyecto integró física eléctrica de manera efectiva, programación en MATLAB y diseño de interfaces gráficas usando AppDesign demostrando como herramientas computacionales pueden modelar fenómenos físicos con aplicación directa en el diagnóstico de enfermedades infecciosas como la malaria.
